\section*{अध्याय \ref{ch:b2:acyl-substitution} --- कार्बोक्सिलिक अम्ल व्युत्पन्न}

\begin{solution}{exo:b2:acyl-substitution:1}
एथेनैमाइड $<$ एथिल एथेनोएट $<$ एथेनोइक ऐनहाइड्राइड $<$ एथेनॉयल क्लोराइड।
\end{solution}

\begin{solution}{exo:b2:acyl-substitution:2}
एथेनोइक अम्ल (+ \ce{HCl}); एथिल एथेनोएट; एथेनैमाइड (+ अमोनियम क्लोराइड, अमोनिया का दूसरा तुल्यांक
\ce{HCl} को पकड़ता है); एथेनोइक ऐनहाइड्राइड (+ \ce{NaCl});
\textit{N},\textit{N}-डाइमेथिलएथेनैमाइड (+ ट्राइएथिलअमोनियम क्लोराइड)।
\end{solution}

\begin{solution}{exo:b2:acyl-substitution:3}
\ce{C6H5COOCH3 + NaOH -> C6H5COONa + CH3OH}। $13.6/136.15 = \qty{0.0999}{mol}$, उतना ही
\ce{NaOH}: \qty{4.00}{g}।
\end{solution}

\begin{solution}{exo:b2:acyl-substitution:4}
$\gamma$-ब्यूटिरोलैक्टोन (ऑक्सोलेन-2-ओन): चार कार्बनों और एक ऑक्सीजन का पंचसदस्यी वलय।
\end{solution}

\begin{solution}{exo:b2:acyl-substitution:5}
अमोनिया कार्बोनिल कार्बन से जुड़ती है; एक प्रोटॉन नाइट्रोजन से ऑक्सीजन (या
निर्गामी समूह) पर जाता है; एथॉक्साइड (प्रोटॉन स्थानांतरण के बाद एथेनॉल के रूप में) निकलता है और \ce{C=O} फिर बनता है:
एथेनैमाइड और एथेनॉल।
\end{solution}

\begin{solution}{exo:b2:acyl-substitution:6}
फ़ीनॉलिक \ce{OH} का ऐसिलन होता है (अम्ल उत्प्रेरण): 2-ऐसीटॉक्सीबेन्ज़ोइक अम्ल; उपोत्पाद एथेनोइक अम्ल।
\end{solution}

\begin{solution}{exo:b2:acyl-substitution:7}
2-मेथिलप्रोपेन-2-ऑल। पहला योग प्रोपेनोन देता है, जो एस्टर से अधिक इलेक्ट्रॉनरागी है और
तुरंत दूसरा मेथिल ग्रहण कर लेता है।
\end{solution}

\begin{solution}{exo:b2:acyl-substitution:8}
ब्यूटेन-1,4-डाइऑल: एस्टर दो प्राथमिक ऐल्कोहॉलों में अपचित होता है, और वलय खुल जाता है।
\end{solution}

\begin{solution}{exo:b2:acyl-substitution:9}
एक तुल्यांक: $x^2/(1 - x)^2 = 4$, $x = 2/3$। तीन तुल्यांक: $x^2/[(1 - x)(3 - x)] = 4$, $3x^2 - 16x
+ 12 = 0$, $x = 0.90$। बनते ही जल को हटाना (डीन--स्टार्क) इसे और आगे बढ़ाता है।
\end{solution}

\begin{solution}{exo:b2:acyl-substitution:10}
\ce{SOCl2} बेन्ज़ॉयल क्लोराइड देता है; डाइएथिलऐमीन के दो तुल्यांकों (या एक और ट्राइएथिलऐमीन) के साथ
यह ऐमाइड देता है। सीधे मिलाए गए अम्ल और ऐमीन अम्ल--क्षारक अभिक्रिया करते हैं: डाइएथिलअमोनियम
बेन्ज़ोएट, जो केवल तीव्र तापन पर ऐमाइड देता है।
\end{solution}

\begin{solution}{exo:b2:acyl-substitution:11}
प्रति ट्राइओलीन तीन KOH: $3 \times 56.11/885.4 = \qty{0.190}{g}$ प्रति ग्राम, अर्थात् \omterm{def:b2:acyl-substitution:saponification}{साबुनीकरण} मान
190।
\end{solution}

\begin{solution}{exo:b2:acyl-substitution:12}
\ce{NaBH4}: केवल कीटोन, द्वितीयक ऐल्कोहॉल तक। \ce{LiAlH4}: कीटोन द्वितीयक
ऐल्कोहॉल तक, एस्टर प्राथमिक ऐल्कोहॉल तक (और ऐल्कोहॉल भाग मुक्त करता है), ऐमाइड ऐमीन तक।
\end{solution}

\begin{solution}{pb:b2:acyl-substitution:1}
\textbf{1.} ओलेइक अम्ल \ce{C18H34O2}; ग्लिसरॉल \ce{C3H8O3}।
\textbf{2.} \ce{C57H104O6}: $57 \times 12.011 + 104 \times 1.008 + 6 \times 15.999 = \qty{885.45}{g/mol}$।
\textbf{3.} तीन।
\textbf{4.} cis द्वि-आबंध शृंखलाओं को मोड़ देता है, जो ठीक से संकुलित नहीं होतीं: दुर्बल अंतराअणुक बल,
निम्न गलनांक।
\textbf{5.} ग्लिसरॉल और तीन सोडियम ओलिएट, अर्थात् साबुन।
\textbf{6.} वसा अम्लों के मेथिल एस्टरों का मिश्रण।
\textbf{7.} \ce{C57H104O6 + 3CH3OH -> C3H8O3 + 3C19H36O2}।
\textbf{8.} मेथेनॉल और क्षारक से बना मेथॉक्साइड आयन मेथेनॉल से कहीं प्रबल नाभिकरागी है;
यह एस्टर कार्बोनिलों पर आक्रमण करता है।
\textbf{9.} एक एस्टर समूह का कार्बोनिल कार्बन, जिस पर \ce{O-}, मेथॉक्साइड का \ce{OCH3} और
ग्लिसरिल ऑक्सीजन हैं: चतुष्फलकीय।
\textbf{10.} \omterm{def:b2:acyl-substitution:transesterification}{पार-एस्टरीकरण} एक साम्य है जिसका स्थिरांक 1 के निकट है; आधिक्य मेथेनॉल उसे
मेथिल एस्टरों की ओर विस्थापित करता है।
\textbf{11.} ग्लिसरॉल निचली परत के रूप में अलग हो जाता है: किसी उत्पाद को हटाना साम्य को आगे बढ़ाता है।
\textbf{12.} नहीं: हर आदान-प्रदान में मेथॉक्साइड का पुनर्जनन होता है।
\textbf{13.} यह क्षारक को उदासीन कर देता है, जिससे सोडियम ओलिएट (साबुन) बनता है: उत्प्रेरक उपभुक्त हो जाता है।
\textbf{14.} यह उनका \omterm{def:b2:acyl-substitution:saponification}{साबुनीकरण} करता है: हाइड्रॉक्साइड (जल और मेथॉक्साइड से) एस्टरों को साबुन और
मेथेनॉल में काट देता है, अनुत्क्रमणीय रूप से।
\textbf{15.} यह क्षारक का उपभोग करता है, मिश्रण को पायसित करता है और ग्लिसरॉल परत को अलग होने से रोकता है।
\textbf{16.} पहले मेथेनॉल के साथ अम्ल-उत्प्रेरित \omterm{def:b2:acyl-substitution:esterification}{एस्टरीकरण} से, जो मुक्त अम्लों को
मेथिल एस्टरों में बदल देता है।
\textbf{17.} जल एस्टरों का जल-अपघटन करता है और उत्प्रेरक को हाइड्रॉक्साइड में बदल देता है, जो \omterm{def:b2:acyl-substitution:saponification}{साबुनीकरण} करता है।
\textbf{18.} $10^6/885.45 = \qty{1129}{mol}$।
\textbf{19.} $3 \times 1129 \times 32.04 = \qty{108.6}{kg}$।
\textbf{20.} $3 \times 1129 \times 296.50 = \qty{1004.6}{kg}$।
\textbf{21.} अंदर $1000 + 108.6 = \qty{1108.6}{kg}$; बाहर $104.0 + 1004.6 = \qty{1108.6}{kg}$।
\textbf{22.} ओलेइक अम्ल के \qty{20}{kg}, $20\,000/282.47 = \qty{70.8}{mol}$, उतना ही \ce{NaOH}:
\qty{2.83}{kg}।
\textbf{23.} $1129 \times 92.09 = \boldsymbol{\approx \qty{104}{kg}}$ ग्लिसरॉल प्रति टन।
\end{solution}
